% Options for packages loaded elsewhere
\PassOptionsToPackage{unicode}{hyperref}
\PassOptionsToPackage{hyphens}{url}
\documentclass[
  12pt,
]{article}
\usepackage{xcolor}
\usepackage{amsmath,amssymb}
\setcounter{secnumdepth}{-\maxdimen} % remove section numbering
\usepackage{iftex}
\ifPDFTeX
  \usepackage[T1]{fontenc}
  \usepackage[utf8]{inputenc}
  \usepackage{textcomp} % provide euro and other symbols
\else % if luatex or xetex
  \usepackage{unicode-math} % this also loads fontspec
  \defaultfontfeatures{Scale=MatchLowercase}
  \defaultfontfeatures[\rmfamily]{Ligatures=TeX,Scale=1}
\fi
\usepackage{lmodern}
\ifPDFTeX\else
  % xetex/luatex font selection
\fi
% Use upquote if available, for straight quotes in verbatim environments
\IfFileExists{upquote.sty}{\usepackage{upquote}}{}
\IfFileExists{microtype.sty}{% use microtype if available
  \usepackage[]{microtype}
  \UseMicrotypeSet[protrusion]{basicmath} % disable protrusion for tt fonts
}{}
\makeatletter
\@ifundefined{KOMAClassName}{% if non-KOMA class
  \IfFileExists{parskip.sty}{%
    \usepackage{parskip}
  }{% else
    \setlength{\parindent}{0pt}
    \setlength{\parskip}{6pt plus 2pt minus 1pt}}
}{% if KOMA class
  \KOMAoptions{parskip=half}}
\makeatother
\usepackage{longtable,booktabs,array}
\usepackage{calc} % for calculating minipage widths
% Correct order of tables after \paragraph or \subparagraph
\usepackage{etoolbox}
\makeatletter
\patchcmd\longtable{\par}{\if@noskipsec\mbox{}\fi\par}{}{}
\makeatother
% Allow footnotes in longtable head/foot
\IfFileExists{footnotehyper.sty}{\usepackage{footnotehyper}}{\usepackage{footnote}}
\makesavenoteenv{longtable}
\setlength{\emergencystretch}{3em} % prevent overfull lines
\providecommand{\tightlist}{%
  \setlength{\itemsep}{0pt}\setlength{\parskip}{0pt}}
\usepackage{bookmark}
\IfFileExists{xurl.sty}{\usepackage{xurl}}{} % add URL line breaks if available
\urlstyle{same}
\hypersetup{
  hidelinks,
  pdfcreator={LaTeX via pandoc}}

\author{SCX}
\date{2026}

\begin{document}

\section{SCX 8-Theorem Paper Mathematical Completeness Audit Report}\label{scx-8-theorem-audit}

\subsection{0. SCX Axiom System Baseline}\label{scx-axiom-baseline}

First clarify the core axioms before auditing:

\begin{itemize}
\tightlist
\item
  \textbf{Theorem 1} (S1): Multi-expert consensus noise detection --- Under A1-A6 assumptions, F1 $\geq$ 1 - (1/$\eta$)$\Sigma$$\rho$\_s$\cdot$exp(-2M$\Delta$\_s$^2$), exponential convergence.
\item
  \textbf{Theorem 2} (S2): Weak feature information-theoretic bound --- F1\_SCX $\leq$ F1\_base + C\_F$\cdot\sqrt{(\delta/2)}$, $\delta$=I($\phi$(X);S).
\item
  \textbf{Theorem 3} (S3): Noise-difficulty unidentifiability --- From observed data, label noise and intrinsic difficulty cannot be distinguished. This is SCX's ``uncertainty principle.''
\item
  \textbf{Theorem 4} (S4): Exact constant minimax optimality --- Bahadur-Rao large deviation exact constant.
\item
  \textbf{Cercis Score}: S(x) = Q(x) + $\eta\cdot$N(x), quality+novelty decomposition.
\item
  \textbf{Situs Operator}: Encodes distributions as metric measure spaces (X, d\_P, $\mu$\_P).
\item
  \textbf{Spring SE-1}: Self-evolution gating, Robbins-Monro convergence guarantee.
\end{itemize}

\begin{center}\rule{0.5\linewidth}{0.5pt}\end{center}

\subsection{1. Per-Paper Audit}\label{per-paper-audit}

\subsubsection{1.1 Theorem 5 --- Active Learning (Cercis Optimal Sampling)}\label{theorem-5-active-learning}

\textbf{Theorem Statement Completeness}: 4/5
- Theorem 5.1 (Optimal Distribution): Complete. Gibbs form p*(x) $\propto$ $\pi$(x)exp($\alpha$S(x)/$\lambda$), with uniqueness proof.
- Theorem 5.2 (Optimal Rate): Statement complete, closed form $\rho$* = $\eta$/(1+$\eta$)$\cdot\Sigma$w\_k$\cdot$V[$\Gamma$\_k]/$\Sigma$w\_k$\cdot\bar{\Gamma}$\_k. But key derivation steps have \textbf{serious logical leaps} (see below).
- Theorem 5.3 (Consistency): Complete. $\sqrt{n}$-consistency, Delta method.

\textbf{Proof Framework Logical Leaps}: 3/5 --- \textbf{Serious problems exist}
- Lemma 1 (Gibbs solution) derivation is standard variational inference, no problem.
- \textbf{Fatal leap in Theorem 5.2 proof lines 414-421}: From $\rho$* = $\eta\cdot\Sigma$w\_k$\bar{\Gamma}$\_k/($\eta\cdot\Sigma$w\_k$\bar{\Gamma}$\_k + $\Sigma$w\_kV[$\Gamma$\_k]) jumping to the final form $\rho$* = $\eta$/(1+$\eta$)$\cdot\Sigma$w\_kV[$\Gamma$\_k]/$\Sigma$w\_k$\bar{\Gamma}$\_k, the paper claims identity $\Sigma$w\_k$\bar{\Gamma}$\_k = ((1+$\eta$)/$\eta$)$\cdot\Sigma$w\_kV[$\Gamma$\_k]. This identity is annotated as ``holds at the variational objective optimum, from equality of marginal benefit and $\eta$-scaled novelty penalty marginal cost.'' \textbf{This is circular reasoning}---using $\rho$*'s definition to derive $\rho$*'s closed form. Algebraically:
- From $\Phi$'($\rho$)=0 we get $\rho$ = $\Sigma$w\_k$\bar{\Gamma}$\_k / ($\Sigma$w\_k$\bar{\Gamma}$\_k + $\eta^{-1}\Sigma$w\_kV[$\Gamma$\_k])
- This form is already a closed form, no additional identity needed
- The paper forces conversion to $\eta$/(1+$\eta$) form by introducing an assumed structural relationship, breaking closure
- \textbf{Assumption 1 (Cercis-Gain Affinity)}: $\bar{\Gamma}$(x) = $\alpha\cdot$S(x) + $\beta$. This is the cornerstone of the entire derivation, but argued only with ``intuition'' and ``Taylor expansion,'' lacking rigorous proof. This is an unverified empirical assumption.

\textbf{Consistency with SCX Axiom System}: 4/5
- Correctly uses Cercis Score S=Q+$\eta$N definition
- $\eta\to0$ recovers pure uncertainty sampling, $\eta\to\infty$ recovers pure diversity sampling, consistent with SCX framework
- But the ``$\eta$/(1+$\eta$)'' saturation factor does not appear in the original SCX framework---this is newly introduced structure

\textbf{Strictly Derivable from Existing Axioms}:
- Theorem 5.1 (Gibbs distribution): Yes, strictly derived from variational principle (no SCX axioms needed)
- Theorem 5.3 (Consistency): Yes, strictly derived from LLN + Delta method

\textbf{Requires New Assumptions}:
- Theorem 5.2 (Sampling rate): No --- needs Assumption 1 (Cercis-Gain Affinity) + additional variational structural relationship
- The risk-adjustment function $\Phi$($\rho$)'s three-term decomposition (expected benefit, marginal diminishing, novelty cost $\times$ variance) is an \textbf{ad hoc construction}, not derived from Cercis axioms

\textbf{Relationship to Other Theorems}: Independent. Referenced by AR-Theorem's detectability result (``T5 optimal active learning strategy''), but AR only uses T5 as a black-box tool.

\begin{center}\rule{0.5\linewidth}{0.5pt}\end{center}

\subsubsection{1.2 Theorem 6 --- Protocol Game (Audit Protocol Adoption)}\label{theorem-6-protocol-game}

\textbf{Theorem Statement Completeness}: 4/5
- Theorem 6.1 (Regret bound): Complete.
- Theorem 6.2 (Nash existence): Complete, standard Glicksberg fixed-point argument.
- Theorem 6.3 (Uniqueness): Complete, explicit $\gamma$\_crit closed form.
- Proposition 6.4 (Distortion): Complete.

\textbf{Proof Framework Logical Leaps}: 4/5
- Existence proof is standard---Glicksberg theorem directly applies to mixed-strategy games with continuous payoffs.
- Uniqueness proof's contraction mapping argument is correct, but one \textbf{leap} exists: the sensitivity of influence function $\phi$\_P to early adopter count is bounded by $\sqrt{2\log N_1/(N_0+1)}$ + $\sigma$\_$\xi$/$\sqrt{N_0}$, and this bound assumes the UCB confidence radius derivative can be approximated by finite differences. This approximation is asymptotically valid for large $N_0$, but requires additional argument for finite samples.
- Lemma ``Influence Monotonicity'' proof relies only on intuitive argument (``increasing n\_P decreases standard error''), lacking formal treatment of UCB index randomness.

\textbf{Consistency with SCX Axiom System}: 3/5
- The paper claims ``Quality(Q) corresponds to protocol quality $\mu$\_P,'' ``Novelty(N) corresponds to UCB exploration bonus.'' This mapping is \textbf{analogical rather than derivational}---the protocol adoption game is not a natural extension of the Cercis framework, but an independently constructed game-theoretic model, then post-hoc analogized to SCX.
- $\gamma$\_crit contains the UCB exploration term $\sqrt{2\log N_1/(N_0+1)}$, which has no direct mathematical connection to Cercis's $\eta$ parameter, only conceptually parallel.

\textbf{Strictly Derivable from Existing Axioms}:
- Theorem 6.1 (Regret): Yes, standard UCB analysis
- Theorem 6.2 (Existence): Yes, standard game theory

\textbf{Requires New Assumptions}:
- Theorem 6.3 (Uniqueness): Requires Gumbel noise assumption + influence sensitivity bound (UCB finite difference approximation)
- The entire model setup (sequential game, Gumbel preference shocks, Gaussian prior) is unrelated to SCX axioms

\textbf{Relationship to Other Theorems}: Indirectly connected to CD-Theorem (information externalities), parallel to TS-Theorem's anchor concept. But \textbf{basically a standalone paper}, attaching to the SCX brand rather than deriving from it.

\begin{center}\rule{0.5\linewidth}{0.5pt}\end{center}

\subsubsection{1.3 Theorem 7 --- Cross-Domain Partition Preservation}\label{theorem-7-cross-domain}

\textbf{Theorem Statement Completeness}: 5/5
- Theorem 7.1 (Main bound): Complete, three components clearly stated.
- Theorem 7.2 (Tightness): Complete, constructs tightness sequence.
- Corollary 7.3 (Optimal granularity): Complete.

\textbf{Proof Framework Logical Leaps}: 4/5
- Five-step proof structure is reasonable.
- \textbf{Step 3 has potential issues}: Claims ``conditional density is L\_k-Lipschitz $\Rightarrow$ entropy functional is L\_k-Lipschitz with respect to Wasserstein metric.'' Entropy is not generally Lipschitz; additional conditions are needed (e.g., log-Sobolev inequality, or bounded conditional density on bounded support). Lemma S2.2 (Entropy Lipschitz Property) gives an argument that H(Y|X=x) is (L/$\sigma$)-Lipschitz, but this assumes bounded conditional variance and near-Gaussian conditional distribution---this is an \textbf{important but unstated assumption}.
- The argument that ``conditioning does not increase average transport cost'' (Step 4) is a correct standard result.

\textbf{Consistency with SCX Axiom System}: 5/5
- Directly starts from Situs conditional mutual information I(Y;P|S)
- Uses Situs encoding and Wasserstein distance, fully within the SCX framework
- Explicitly cites T7 cross-domain transfer theorem

\textbf{Strictly Derivable from Existing Axioms}:
- Main bound: partially derivable (needs Assumption 1 + entropy Lipschitz condition)
- Asymptotic tightness: Yes, constructive proof

\textbf{Requires New Assumptions}:
- Assumption 1 (Bounded Conditional Densities with Lipschitz property)
- Lemma S2.2's entropy Lipschitz argument requires Gaussianity or equivalent condition

\textbf{Relationship to Other Theorems}: Cited by CD-Theorem's causal relay theorem as ``T7 cross-domain fidelity.'' Among all 8 papers, \textbf{mathematically the most self-consistent}.

\begin{center}\rule{0.5\linewidth}{0.5pt}\end{center}

\subsubsection{1.4 CD-Theorem --- Causal Discovery (Noise vs Confounding)}\label{cd-theorem}

\textbf{Theorem Statement Completeness}: 4/5
- CD-Theorem 1 (Strong Inseparability): Complete.
- CD-Theorem 2 (Weak Separability): Statement complete but \textbf{extremely conditional}.
- CD-Theorem 3 (Causal Relay): Complete.

\textbf{Proof Framework Logical Leaps}: 3/5
- \textbf{Strong Inseparability proof leap}: Step 1 claims S(x) = S$_0$(x) + $\eta\cdot\delta\_\varepsilon$(x) + $\beta\cdot\delta$\_z(x), then asserts that when $\eta\perp$z, the joint distribution factorization makes $\delta$\_z indistinguishable from effective noise. This is a direct analogical generalization of Theorem 3, but the \textbf{argument is overly brief}---need to prove that the joint distribution of (S, $\nabla$S, H\_S) is indeed identical in both worlds, not just the marginal distribution of S. Theorem 3's original proof explicitly constructs joint distribution equivalence for the two worlds; CD-Theorem does not do this explicit construction.
- \textbf{Curl statistic problem}: $\tau$ = (1/|X|)$\int\|\nabla\times\hat{\nabla}S\|^2$dx. This statistic computationally relies heavily on the curl definition on the Situs manifold. The paper does not give a concrete construction of the curl operator on the Situs manifold, only vaguely references ``Situs manifold in the SCX framework.'' Geometrically, curl requires a connection structure, while Situs encoding only provides a metric---\textbf{from metric to curl requires additional mathematical structure} (Levi-Civita connection may not exist on general metric spaces).
- \textbf{Causal Relay}: Chain derivation is reasonable, but the interaction term C$\cdot\varepsilon\cdot\delta$ in g($\varepsilon$,$\delta$) does not give C's definition or bounding method.

\textbf{Consistency with SCX Axiom System}: 4/5
- Directly extends Theorem 3 to causal discovery setting
- Situs skeleton preservation assumption is a major extension of the Situs operator, but the paper honestly labels it as an open problem
- Causal Relay correctly cites T7

\textbf{Strictly Derivable from Existing Axioms}:
- Strong Inseparability: Claimed as direct corollary of Theorem 3, but proof is overly brief

\textbf{Requires New Assumptions}:
- Situs skeleton preservation assumption (core open problem)
- Existence of curl operator on Situs manifold
- All three theorems have honest annotations ([Rigorous]/[Conditionally Rigorous]/[Open])

\begin{center}\rule{0.5\linewidth}{0.5pt}\end{center}

\subsubsection{1.5 FA-Theorem --- Federated Audit (ZK Multi-Party)}\label{fa-theorem}

\textbf{Theorem Statement Completeness}: 4/5
- FA-Theorem 1 ($\varepsilon$-Equivalence): Complete, four error sources.
- FA-Theorem 2 (Info Lower Bound): Complete.
- ZK construction sketch: Complete but only an outline.

\textbf{Proof Framework Logical Leaps}: 3/5
- \textbf{Core unsolved problem}: Existence of Situs homomorphic composition operator $\oplus$. Without it, the $\varepsilon$-Equivalence theorem has only conditional meaning. The paper honestly admits this is an open problem.
- \textbf{I-MMSE gap derivation} (Step 2 of Lower Bound): Uses Guo-Shamai-Verd\'{u} I-MMSE relation, but this relation applies to Gaussian channels, while the Cercis Score here is non-Gaussian. The argument requires generalization or additional justification.
- \textbf{``Worst construction''} (Step 3): Claims ``adversary can adjust \~{S}\_k to maximize bias,'' but does not give the specific form of the adjustment strategy. Needs to prove existence of a tampering that makes the deviation between E[T\_fed] and E[T\_full] reach the claimed lower bound.

\textbf{Consistency with SCX Axiom System}: 4/5
- Correctly uses Cercis Score and Situs encoding
- Audit Sword principle is an extension of the SCX framework
- Lower bound directly derived from Theorem 3

\textbf{Strictly Derivable from Existing Axioms}:
- Info Lower Bound (Theorem 2): Yes, derived from Theorem 3 + data processing inequality

\textbf{Requires New Assumptions}:
- Situs homomorphic composition operator $\oplus$ (core open problem)
- Feasibility of ZK circuit construction (cryptographic engineering problem, not mathematical)

\textbf{Relationship to Other Theorems}: Cited by CD-Theorem and AR-Theorem as ``Audit Sword'' principle's federated extension. \textbf{The most ambitious extension paper}, but also the most dependent on unsolved mathematical problems (Situs homomorphism).

\begin{center}\rule{0.5\linewidth}{0.5pt}\end{center}

\subsubsection{1.6 AR-Theorem --- Adversarial Robustness (Cercis Taxonomy)}\label{ar-theorem}

\textbf{Theorem Statement Completeness}: 4/5
- AR-Theorem 1 (Reduction): Complete, conditions clear.
- AR-Theorem 2 (Detectability): Complete, sample complexity lower bound.
- Conjecture (Phase Transition): Statement complete but only a conjecture.

\textbf{Proof Framework Logical Leaps}: 3/5
- \textbf{Reduction theorem key leap}: Proof claims from $\|\nabla S\|$/S $\leq$ $\eta$ (``from Cercis operator definition'') that A(x) $\leq$ $\eta\cdot\varepsilon$. This inequality \textbf{is not a standard result in the SCX framework}---Cercis operator defines S=Q+$\eta$N, which does not directly constrain $\|\nabla S\|$/S. This requires the additional assumption that $\nabla$Q and $\nabla$N are bounded by Q and N proportionally.
- \textbf{Hessian bound in Taylor expansion}: $\|$H\_S$\|\_\text{op}$ $\leq$ L(x) is correct (Lipschitz gradient $\Rightarrow$ Hessian bounded), but only when using spectral norm and the function is twice differentiable.
- \textbf{Detectability theorem's lower bound}: Uses Le Cam two-point method to derive $\chi^2$ divergence = $\Theta$($\gamma^2$). This derivation is \textbf{extremely brief}---needs explicit computation of the $\chi^2$ divergence between gradient field distributions under H$_0$ and H$_1$, which depends on specific assumptions about A(x)'s distribution. The paper does not give this computation.
- \textbf{Upper bound (T5 achievability)}: Claims Theorem 5 provides matching upper bound, but T5's optimality is about maximizing I(S;query), not about rate optimality for detecting adversarial structure. This is \textbf{different meaning of optimality}---T5 optimizes sample efficiency for label acquisition, AR needs sample efficiency for detecting anomalous gradient structure.

\textbf{Consistency with SCX Axiom System}: 3/5
- Ternary decomposition S=Q+$\eta$N+$\gamma$A is a major extension of the original Cercis binary decomposition
- Whether introducing a third primitive (A) within the SCX framework is an \textbf{axiomatic-level question}
- The gradient-field diagnostic D\_adv construction is novel but lacks clear definition within the SCX framework

\textbf{Strictly Derivable from Existing Axioms}:
- Reduction's Lipschitz sufficient condition: derivable, but uses an inequality ($\|\nabla S\|$/S $\leq$ $\eta$) not in SCX axioms

\textbf{Requires New Assumptions}:
- Plausibility of ternary decomposition S=Q+$\eta$N+$\gamma$A
- $\|\nabla S\|$/S $\leq$ $\eta$ (Cercis gradient-value ratio bounded)
- Incomplete correlation of adversarial exposure A with novelty N (I(A;N|S)>0)
- Curl operator on Situs manifold (shared with CD-Theorem)

\textbf{Relationship to Other Theorems}: Explicitly cites CD (adversarial signals confused with causal signals), HC (humans detect adversarial samples), TS (S\_crit drift). But among the 8 papers, \textbf{mathematically weakest}---two theorems heavily depend on unverified assumptions, core conjecture is completely open.

\begin{center}\rule{0.5\linewidth}{0.5pt}\end{center}

\subsubsection{1.7 TS-Theorem --- Temporal SCX (Drift \& Self-Audit)}\label{ts-theorem}

\textbf{Theorem Statement Completeness}: 4/5
- TS-Theorem 1 (Drift Detectability): Complete, sample size formula.
- TS-Theorem 2 (Drift Attribution): Complete, split into anchor-free/anchored cases.
- TS-Theorem 3 (Sprint Self-Audit Limit): Complete, Lyapunov condition.

\textbf{Proof Framework Logical Leaps}: 4/5
- \textbf{Drift Detectability proof is reasonable}, but has one implicit assumption: Lipschitz relationship between Situs distance and $\eta$ drift $|\eta_{t+1} - \eta_t| \leq \kappa\cdot$d\_Situs + $|\xi\_t|$. This relationship is reasonably derivable from T7's Lipschitz property.
- \textbf{Variance formula (Equation 14)} contains $\sigma^2$\_est term without defining or deriving it. Under H$_0$, $\eta$'s estimation variance should be derivable from Cercis Score distribution, but the paper directly asserts a form.
- \textbf{Attribution unidentifiability construction} (Config A vs Config B) is persuasive, but only for two configurations; does not prove \textbf{all} possible decompositions are unidentifiable. Needs more general unidentifiability proof.
- \textbf{Lyapunov convergence theorem} is standard Foster-Lyapunov argument, no problem. Divergence condition derivation is also natural.

\textbf{Consistency with SCX Axiom System}: 4/5
- Explicitly temporalizes T7 cross-domain transfer theorem (Situs distance becomes Situs distance between time slices)
- Sprint self-audit problem is a natural temporal extension of the SCX framework
- Anchor concept parallels HC-Theorem's anchor knowledge

\textbf{Strictly Derivable from Existing Axioms}:
- Drift Detectability: Yes, from T7 + standard hypothesis testing
- Attribution unidentifiability: Yes, via counterexample construction
- Anchored estimation: Yes, directly derived from anchor definition

\textbf{Requires New Assumptions}:
- Lipschitz constant $\kappa$ from Situs distance to $\eta$
- Existence of anchor set (practical but not theoretically required)
- Contraction rate $\rho$ in Lyapunov condition (needs derivation from Sprint update rule)

\begin{center}\rule{0.5\linewidth}{0.5pt}\end{center}

\subsubsection{1.8 HC-Theorem --- Human-AI Collaborative Audit}\label{hc-theorem}

\textbf{Theorem Statement Completeness}: 5/5
- HC-Theorem 1 (Collaborative Gain): Complete.
- HC-Theorem 2 (Budget Decay): Complete, $\sqrt{B_H/n}$ scaling law.
- HC-Theorem 3 (Absolute Boundary): Complete.

\textbf{Proof Framework Logical Leaps}: 4/5
- \textbf{Collaborative Gain logic}: ``If human judgment carries information beyond SCX, then collaborative audit strictly dominates pure AI.'' This is a conditional implication, logically correct, but \textbf{the condition itself (I(J\_H;$\varepsilon$|S)>0) is an empirical hypothesis}. The paper honestly admits this is an open problem.
- \textbf{Budget Decay derivation has subtle issues}: Step 1 claims $\Omega$\_collab = $\Omega$\_AI + (B\_H/n)$\cdot\Delta$\_I$\cdot\eta$\_eff, but Step 2 via Cram\'{e}r-Rao argument gives $\eta$\_eff $\leq$ c/$\sqrt{B_H/n}\cdot$1/$\Delta$\_I, yielding $\sqrt{B_H/n}$ scaling. This compound derivation needs to prove: optimal selection based on Cercis Score indeed produces $\sqrt{B_H/n}$ efficiency loss. The paper's argument is overly brief---``optimal selection introduces inter-sample dependence, reducing effective sample size from B\_H to $\sqrt{B_H\cdot n}$'' needs more rigorous treatment.
- \textbf{Absolute Boundary theorem}: This is a direct logical extension of Theorem 3, argument is solid. ``Perfect noise'' concept is clear: I($\varepsilon$; all observables | X) = 0 $\Rightarrow$ absolutely indistinguishable.

\textbf{Consistency with SCX Axiom System}: 5/5
- Directly extends Theorem 3 to any cognitive system
- ``Multi-lens'' perspective is a philosophical deepening of the SCX framework
- Establishes a continuum from Theorem 3 (AI-lens) to any lens

\textbf{Strictly Derivable from Existing Axioms}:
- Absolute Boundary (Theorem 3): Yes, direct corollary of Theorem 3
- Budget Decay: Derivable from Theorem 3 + information theory, but $\sqrt{B_H/n}$ scaling's rigor needs more work

\textbf{Requires New Assumptions}:
- Information asymmetry condition I(J\_H;$\varepsilon$|S)>0 (empirical hypothesis)
- Calibration condition (human uncertainty monotonic in Cercis Score)
- Efficiency factor $\eta$\_eff's upper bound derivation needs stronger regularity conditions

\textbf{Relationship to Other Theorems}: Connected to CD (intersection of human cognition and SCX signals), AR (humans detect adversarial samples), TS (anchors). Among the 8 papers, \textbf{philosophically deepest, with the most thorough understanding of Theorem 3}.

\begin{center}\rule{0.5\linewidth}{0.5pt}\end{center}

\subsection{2. Cross-Paper Analysis}\label{cross-paper-analysis}

\subsubsection{2.1 Theorem Duplication or Contradiction}\label{theorem-duplication}

\begin{itemize}
\tightlist
\item
  \textbf{Duplication}: No direct duplication. CD-Theorem's Strong Inseparability overlaps in spirit with Theorem 3 but addresses different questions (noise vs difficulty vs noise vs confounding).
\item
  \textbf{Potential contradictions}:

  \begin{itemize}
  \tightlist
  \item
    AR-Theorem's Detectability claims to break Theorem 3 at the $\nabla$S level (``Theorem 3's barrier applies at the S level, not the $\nabla$S level''), while Theorem 3's original statement covers ``Cercis Score and its derivatives.'' If Theorem 3's original proof indeed covers $\nabla$S, then AR's claim contradicts Theorem 3. \textbf{This requires checking whether Theorem 3's proof truly covers gradient information.} Theorem 3's original text says ``observables accessible to an SCX auditor: the Cercis Score, its gradient, its Hessian,'' concluding ``No algorithm operating on these observables can tell the worlds apart.'' AR-Theorem claims $\nabla$S contains higher-order structure (curl) not exhausted by Theorem 3. \textbf{This is a potential internal contradiction---either Theorem 3's claim is weakened by AR, or AR's claim is invalid.}
  \item
    T5's $\eta$/(1+$\eta$) factor's semantics is not fully consistent with the original SCX $\eta$ parameter: in the original framework $\eta$ is a continuous exploration-exploitation tradeoff, while in T5 $\eta$ affects sampling rate through the $\eta$/(1+$\eta$) saturation factor. Should these two $\eta$'s be considered the same parameter?
  \end{itemize}
\end{itemize}

\subsubsection{2.2 Weakest Paper: \textbf{AR-Theorem (Adversarial Robustness)}}\label{weakest-paper}

\begin{itemize}
\tightlist
\item
  \textbf{Reasons}:

  \begin{enumerate}
  \def\labelenumi{\arabic{enumi}.}
  \tightlist
  \item
    One core result (Reduction) depends on inequality $\|\nabla S\|$/S $\leq$ $\eta$, which is not in the SCX axioms
  \item
    Detectability's $\chi^2$ divergence computation is not explicitly given
  \item
    Core problem (Phase Transition Conjecture) is entirely an open conjecture
  \item
    Ternary decomposition S=Q+$\eta$N+$\gamma$A is a major modification of SCX axioms, but the paper does not argue for the legitimacy of this modification
  \item
    Potential contradiction with Theorem 3
  \item
    Le Cam lower bound derivation is overly brief, may contain hidden assumptions
  \end{enumerate}
\item
  \textbf{Honest roast}: This paper has a good idea (adversarial samples may constitute a third category), but the mathematical execution is insufficient. The conditionality of the two ``theorems'' is too strong, and core conclusions are mostly ``if... then...'' conditional statements, where the conditions themselves are unmet or unverified. Essentially a \textbf{positioning paper + research program}, not a resolved theorem.
\end{itemize}

\subsubsection{2.3 Strongest Paper: \textbf{Theorem 7 (Cross-Domain Partition Preservation)}}\label{strongest-paper}

\begin{itemize}
\tightlist
\item
  \textbf{Reasons}:

  \begin{enumerate}
  \def\labelenumi{\arabic{enumi}.}
  \tightlist
  \item
    Most complete theorem statements, every symbol has a clear definition
  \item
    Five-step proof decomposition is clear, with well-developed lemmas
  \item
    Constructs asymptotic tightness, the highest standard of theoretical proof
  \item
    Highest consistency with SCX framework---directly from Situs conditional mutual information
  \item
    Component decomposition (geometric $\times$ statistical $\times$ irreducible) provides actionable practical guidance
  \item
    Extensions (p-Wasserstein, multi-target domain, adaptive refinement) demonstrate the theory's richness
  \item
    \textbf{The only significant weakness} (entropy Lipschitz requiring Gaussian assumption) is discussed in the appendix
  \end{enumerate}
\item
  \textbf{Runner-up candidate}: HC-Theorem --- the philosophical deepening of Theorem 3 is outstanding, three-layer structure (AI-only $\to$ human-augmented $\to$ absolute limit) has perfect logic. Budget Decay's $\sqrt{B_H/n}$ scaling is a beautiful information-theoretic result. The only weakness is that the core hypothesis depends on empirical validation.
\end{itemize}

\begin{center}\rule{0.5\linewidth}{0.5pt}\end{center}

\subsection{3. Global Evaluation Matrix}\label{global-evaluation-matrix}

\begin{longtable}[]{@{}llllll@{}}
\toprule\noalign{}
Paper & Completeness & Logical Rigor & SCX Consistency & New Assumption Dependence & Overall Score \\
\midrule\noalign{}
\endhead
\bottomrule\noalign{}
\endlastfoot
T5-Active & 4/5 & 3/5 & 4/5 & High & \textbf{B+} \\
T6-Protocol & 4/5 & 4/5 & 3/5 & High & \textbf{B} \\
T7-CrossDomain & 5/5 & 4/5 & 5/5 & Medium & \textbf{A} \\
CD-Causal & 4/5 & 3/5 & 4/5 & High & \textbf{B-} \\
FA-Federated & 4/5 & 3/5 & 4/5 & Very High & \textbf{B-} \\
AR-Adversarial & 4/5 & 2/5 & 3/5 & Very High & \textbf{C+} \\
TS-Temporal & 4/5 & 4/5 & 4/5 & Medium & \textbf{B+} \\
HC-Human & 5/5 & 4/5 & 5/5 & Medium & \textbf{A-} \\
\end{longtable}

\begin{center}\rule{0.5\linewidth}{0.5pt}\end{center}

\subsection{4. Key Findings Summary}\label{key-findings-summary}

\begin{enumerate}
\def\labelenumi{\arabic{enumi}.}
\item
  \textbf{T5's Theorem 5.2 proof has an algebraic/logical leap}---the final step of the closed form introduces circular reasoning. But this problem is fixable (directly use $\rho$ = $\Sigma\bar{\Gamma}$/($\Sigma\bar{\Gamma}$+$\eta^{-1}\Sigma$V) as the closed form, no need to force conversion to $\eta$/(1+$\eta$) form).
\item
  \textbf{AR has a potential contradiction with Theorem 3}---AR claims the gradient field contains information not covered by Theorem 3, but Theorem 3 explicitly states it covers gradients. Needs CLARIFICATION.
\item
  \textbf{Situs homomorphic composition operator ($\oplus$) is FA-Theorem's core bottleneck and also the blocking point for the entire federated SCX program}. Without $\oplus$, FA-Theorem degrades to a conditional statement.
\item
  \textbf{CD-Theorem's curl statistic depends on the connection structure on the Situs manifold}, which does not naturally exist in pure metric spaces. This is a geometric problem shared by CD and AR.
\item
  \textbf{Among the 8 papers, 4 have core results that are to some extent ``conditional theorems''} (if X holds, then Y holds), where X is often an open problem or empirical hypothesis. This is not a bad thing---it correctly positions the frontier of SCX research---but readers need to know which are closed and which are open.
\item
  \textbf{The strongest mathematics is in T7} (information-theoretic bound + tightness construction), \textbf{the deepest philosophy is in HC} (multi-lens theory + perfect noise absolute boundary), \textbf{the most practical is in T5} (closed-form sampling rate), \textbf{the boldest are in FA and AR} (but with the weakest mathematical foundations).
\end{enumerate}

\end{document}
